\section*{Chapter \ref{ch:fm:hiring-and-compensation} --- Hiring and Compensation}

\begin{solution}{exo:fm:hiring-and-compensation:1}
$1-\Phi(1.2)=11.5\%$ (about one year in nine) and $\Phi(-1.0)=15.9\%$ (one in six).
\end{solution}

\begin{solution}{exo:fm:hiring-and-compensation:2}
$0.2\times(120-60-10)=10$ and $0.06\times120=7.2$.
\end{solution}

\begin{solution}{exo:fm:hiring-and-compensation:3}
The year's cash, 0.6, plus a quarter of the deferred part of each of the four previous awards, $4\times0.1$: a full award of 1.
\end{solution}

\begin{solution}{exo:fm:hiring-and-compensation:4}
$s=4$: luck share $100/116=86.2\%$, $T=100/(16\times0.5625)=11.1$ years. $s=8$: 61.0\% and 2.78 years.
\end{solution}

\begin{solution}{exo:fm:hiring-and-compensation:5}
Above 100\% it needs shareholders' approval (up to 200\%, by a qualified majority). At least 40\% must be deferred over four to five years (60\% if
the amount is particularly high) and at least 50\% paid in shares or equivalent instruments.
\end{solution}

\begin{solution}{exo:fm:hiring-and-compensation:6}
Pay is $r\max(P\&L,0)$, convex: more variance raises its expectation while losses cost the trader nothing below zero. Limits cap the risk; deferral
with \omterm{def:fm:hiring-and-compensation:malus}{malus} makes later losses reduce earlier pay.
\end{solution}

\begin{solution}{exo:fm:hiring-and-compensation:7}
About 0.76, against 0.64 after five years and a closed-form 0.78 for average P\&L.
\end{solution}

\begin{solution}{exo:fm:hiring-and-compensation:8}
One year of \$25 million is weak evidence of skill (luck is 86\% of a year's variance on this desk); a guarantee pays whatever the next years
show and removes the link between pay and results; and where EU rules apply guaranteed variable pay is allowed only on hiring. Offer deferred pay
whose value depends on future results.
\end{solution}

\begin{solution}{pb:fm:hiring-and-compensation:1}
\begin{enumerate}
\item See \cref{def:fm:hiring-and-compensation:structured}; it is the top-ranked selection procedure in the 2022 revision of the validity
research.
\item A timed problem set, a structured phone interview, an on-site loop of \omterm{def:fm:hiring-and-compensation:structured}{structured interviews} scored independently, and a decision meeting on
the scores: cheapest valid tests first.
\item Most candidates will not be hired, so a late-stage test spends its effort on people an early test could have removed.
\item See \cref{def:fm:hiring-and-compensation:pool}.
\item 6 at revenue 100; 0 at revenue 70.
\item It pays nothing in a year the desk does not earn its cost of capital; a revenue pool still pays.
\item See \cref{prop:fm:hiring-and-compensation:years}.
\item 86\% and 11.1 years.
\item 87\%.
\item One year: 0.35, 0.34, 0.37; five years: 0.64, 0.63, 0.61 (formulaic, risk-adjusted, discretionary).
\item Judgement on the P\&L is a shrinkage of it and cannot add information; observations of decision quality, risk and contribution to others can.
\item About 0.76.
\item See \cref{def:fm:hiring-and-compensation:deferral,def:fm:hiring-and-compensation:malus,def:fm:hiring-and-compensation:buyout}.
\item One full year's award.
\item The unvested balance, one year's award, in its own deferred instruments on a similar schedule.
\item Variable pay at most 100\% of fixed (200\% with approval), at least 40\% deferred over four to five years, at least 50\% in instruments; the UK
removed the ratio limit from 31 October 2023.
\item \omterm{def:fm:hiring-and-compensation:malus}{Malus} reduces what the firm still holds; \omterm{def:fm:hiring-and-compensation:malus}{clawback} must recover money already paid.
\item Not with a guarantee: one year is mostly luck; match with deferred pay tied to future results, if at all.
\item 87\% of the variance of a year's formulaic pay is luck; about twelve years of P\&L (11.1 in closed form) for a correlation of 0.8 with skill.
\item Pay a share of profit after the \omterm{def:fm:limits-and-capital-allocation:raroc}{capital charge}, allocated by results averaged over several years and by observed decision quality; defer a
large part with \omterm{def:fm:hiring-and-compensation:malus}{malus}, so that later results adjust earlier pay.
\end{enumerate}
\end{solution}

\begin{solution}{iq:fm:hiring-and-compensation:1}
To expose pay to losses that appear later (\omterm{def:fm:hiring-and-compensation:malus}{malus}), to align it with longer-term results, and to retain.
\iqlookfor{\omterm{def:fm:hiring-and-compensation:malus}{malus} and retention.}
\end{solution}

\begin{solution}{iq:fm:hiring-and-compensation:2}
\omterm{def:fm:hiring-and-compensation:malus}{Malus} reduces unvested awards; \omterm{def:fm:hiring-and-compensation:malus}{clawback} recovers pay already vested or paid.
\iqlookfor{the vesting boundary.}
\end{solution}

\begin{solution}{iq:fm:hiring-and-compensation:3}
Correlation $\sqrt{0.16/(0.16+1/T)}$; 0.8 needs $T=11.1$ years.
\iqlookfor{signal-to-noise and the $1/T$ term.}
\end{solution}

\begin{solution}{iq:fm:hiring-and-compensation:4}
Compare expected pay over several years, including the unvested balance forfeited, the vesting and \omterm{def:fm:hiring-and-compensation:malus}{malus} terms, and the capital and limits each firm
gives you.
\iqlookfor{the unvested balance and risk-adjusted expectations.}
\end{solution}

\begin{solution}{iq:fm:hiring-and-compensation:5}
Pay is convex in P\&L; control it with risk limits, \omterm{def:fm:limits-and-capital-allocation:raroc}{capital charges} and deferral with \omterm{def:fm:hiring-and-compensation:malus}{malus}.
\iqlookfor{convexity and its antidotes.}
\end{solution}

\begin{solution}{iq:fm:hiring-and-compensation:6}
Add observations with a higher signal-to-noise ratio than P\&L (decision quality, forecast calibration, execution quality), shrink P\&L towards
them, and defer so that the P\&L's later evidence can adjust the award.
\iqlookfor{more information, shrinkage and deferral.}
\end{solution}
